\chapter*{Conferencia II}
\addcontentsline{toc}{chapter}{LECTURA I - El timbre y el carácter del violín}
Señores:

Estaba trabajando en un violín usado,
un violín de defectos de tono "ruidosos", un violín perteneciente
Para un jugador "comercial". La regradación estaba completa, una nueva barra ajustada, el trabajo de revestimiento interior terminado y todo listo para el montaje cuando el dueño me llamó para que me diera prisa. Darme prisa significaba trabajar de noche, y para complacer a mi amigo, decidí hacer el trabajo de montaje esa misma noche. En consecuencia, regresé a mi taller, encendí una lámpara de 20 candelas colocada frente a un reflector y me senté justo frente a la lámpara.
Extendiendo la mano hacia mi derecha, tomé la caja de resonancia terminada. Al colocarla entre mi vista y la lámpara, noté algunas manchas oscuras en la superficie interior. Pensando que de alguna manera la había ensuciado, dejé la caja de resonancia sobre el banco con la intención de limpiarlas con un paño. Pero, al dejarla sobre el banco, no se veían manchas. Habían desaparecido. La superficie interior estaba limpia y brillante. Pero, sin duda, vi manchas oscuras. Volví a acercar la caja de resonancia a la lámpara. Ahí estaban, claramente visibles; varias; algunas grandes, otras pequeñas. Debían de ser manchas en la superficie barnizada. Examiné el barniz con detenimiento. No aparecían manchas. Posiblemente estas manchas se debieran a opacidades en la madera, pero no encontré ninguna; la madera estaba limpia y brillante. Volví a acercar la caja de resonancia a la lámpara. Había una, dos, tres, seis zonas opacas, ubicadas en la zona superior de la caja de resonancia, la que produce el sonido.
Inmóvil, contemplo esas zonas nubosas, pensando,
Pensando mucho en una explicación. Lentamente llegó la explicación, lentamente por supuesto. No soy brillante; nunca me he hecho pasar por un prodigio; Soy lento, pero soy persistente, y también fumador. Como saben, fumar es un requisito indispensable para el filósofo. No podemos separar la pipa del alemán, ni el mundo puede igualarlo en la resolución de problemas complejos. Imitando al filósofo, enciendo mi pipa. De inmediato, a través de nubes de humo, veo la causa de estas zonas nubosas. Así es como el humo ayuda a la filosofía. El gran Hahnemann estableció el principio de que "lo semejante cura lo semejante"; o, como Hahnemann afirma en un latín exquisito, "Similia similibus curantur". Creo en Hahnemann en lo que respecta al humo; sin embargo, ese gran pensador logró traspasar la dura piel de los médicos convencionales con el hecho punzante de que nuestras dosis a menudo exceden la generosidad. Y ahora yo mismo, uno de los médicos convencionales, voy a demostrar que lo infinitesimal también puede causar "ruido" dentro del violín. Creo en el azar. Con el tiempo, les presentaré dos sucesos accidentales que apuntan a... directamente a soluciones de problemas de sonido de violín hasta ahora considerados irresolubles. Por lo tanto, cuando me enfrento a problemas irresolubles, rezo por la casualidad.
Solo por llevar accidentalmente esta caja de resonancia
Con su superficie convexa hacia la lámpara, aún podría estar buscando la causa del "ruido" en el tono del violín; es más, podría no haber encontrado nunca la causa. En la siguiente narración se puede ver claramente la evidencia incuestionable que se ofrece 
Por casualidad, pues, a la luz de las teorías del sonido musical, esta evidencia se presenta de forma muy clara.
Para aclarar mi razonamiento, recurro a algunos hechos conocidos de la filosofía natural: Así, los rayos de luz que inciden sobre la superficie convexa de un cuerpo cóncavo-convexo (como la caja de resonancia del violín), tras atravesarlo, se desvían de una línea recta al salir de la superficie cóncava y, desde allí, convergen hacia un foco; por lo tanto, los objetos sobre la superficie cóncava se magnifican; pero, invirtiendo la dirección de propagación, los rayos de luz que salen de la superficie convexa se dispersan; por consiguiente, los objetos sobre la superficie convexa parecen disminuir de tamaño; por lo tanto, si hubiera llevado esta caja de resonancia con su superficie convexa hacia la lámpara, no habría podido ver esas zonas borrosas. [El barniz transparente sobre la superficie convexa de esta caja de resonancia facilita enormemente la transmisión de los rayos de luz; y, al aplicar esta prueba para un trabajo de graduación perfecto sobre madera de caja de resonancia "en bruto", es necesario primero cubrir la superficie convexa con algún medio transparente, como aceite, barniz transparente o, mejor aún, una mezcla de estas dos sustancias disponibles.]
Al realizar el trabajo de graduación en esta caja de resonancia, supuse que estaba haciendo un buen trabajo; usando los calibradores con cuidado; sin embargo, no logré hacer un trabajo perfecto, como se demostrará ahora. Si esta caja de resonancia hubiera sido reemplazada sin descubrir esas áreas nubosas, aún se habrían mezclado con su tono ondas sonoras con un tono determinado.
teclas disonantes o "ruido". Como antes, sosteniendo la caja de resonancia a contraluz, con un lápiz delimito una de esas áreas y, con una espátula, retiro un poco de madera. Ahora, en esta área, la luz pasa con mayor libertad. Continúo así con la espátula hasta que esta área delimitada tenga la misma luminosidad que las áreas adyacentes.
"Qoctor, ¿cuánta madera quitó?" No puedo responder a esta pregunta con precisión; pero diga que tal vez la centésima parte de una pulgada en algunos lugares; en otros lugares de mayor opacidad, una cantidad mayor. Debe entender que la forma de graduación aplicada a esta caja de resonancia es de 9,64 pulgadas en la posición del puente; de ahí, en dirección ascendente o descendente, hasta 3,32 pulgadas en el borde f lo más cerca posible del trabajo. Por lo tanto, la precisión absoluta en la disminución de los espesores de la caja de resonancia se convierte en una cuestión de dificultad; de hecho, esta forma de graduación de la caja de resonancia no es superada en dificultad por ninguna otra forma conocida. También debe entender que donde el espesor es igual a 9,64, pero pasa poca o ninguna luz; también, que la mayor luminosidad, en esta placa, ocurre en el espesor de 3,32 pulgadas; que entre estos dos puntos, la luminosidad disminuye o aumenta con una regularidad exactamente proporcional a la disminución o aumento del espesor de la caja de resonancia; Por lo tanto, las irregularidades en el grosor producen zonas nubosas al interferir con el paso de los rayos de luz. Recuerdo las reglas que rigen la afinación. Menciono que las irregularidades de una centésima de pulgada en la lengüeta del órgano, en el tubo del órgano o en un diapasón son inadmisibles. \label{sobretonos_dis}
¿Es tal irregularidad también inadmisible en la caja de resonancia del violín? ¿He encontrado, por fin, la causa del "ruido" en el sonido del violín? Me hormiguean los nervios. Llevo cuarenta años trabajando con violines usados tratando de encontrar una fórmula para erradicarlo.
"Ruido" proveniente de ello; funcionó con la cantidad habitual de fallos "ruidosos". Más rápido de lo que puedo percibir, la información llegó a mis sentidos: esas leves irregularidades en el grosor de la caja de resonancia podían producir un tono estridente, penetrante, desgarrador e inconmensurablemente agudo; un tono de ese tipo por cada irregularidad en el área productora de sonido de la caja de resonancia del violín.
Habrás oído hablar de cómo los "buscadores" pueden "buscar durante toda la vida sin encontrar una pista"; y cómo, cuando encuentran una pista, inmediatamente celebran el evento volviéndose locos.
Señores: El sonido de este violín no está acompañado de ondas sonoras con tonos disonantes; y ahora conocen la causa del "ruido" en el sonido del violín. En relación con la palabra "violín", la palabra "ruido" evoca vívidamente un sufrimiento incalculable; por lo tanto, no insistiré más en este tema de lo necesario para que comprendan completamente esos armónicos agudos y disonantes que provienen de la imperfecta caja de resonancia del violín. La ciencia afirma que el sonido musical comprende, en cuanto a tono, de 27 a 4000 vibraciones por segundo; y, dado que el "ruido" del violín es exasperantemente agudo, se deduce que posee un tono que supera las 4000 vibraciones por segundo.
La ciencia también afirma que el sonido puede tener un tono determinado.
tan alto que hace imposible su medición precisa. Así, esos agudos tonos de violín (no armónicos legítimos) se convierten en objetos legítimos para "adivinar" su tono; y cuando adivinar se vuelve legítimo, es culpa del "adivinador" si su "adivinación" es demasiado baja; por lo tanto, "adiviné" que las vibraciones de estos "ruidos" que sugieren suicidio oscilan entre 20 000 y 20 000 000 vibraciones por segundo. Al hacer esta "adivinación", admito que mis cifras pueden ser grandes en proporción al tamaño de mi oído; de todos modos, la oportunidad de exponer estas pequeñas y molestas molestias a su verdadera luz resulta muy reconfortante para mis nervios, que han sufrido durante mucho tiempo.
Dicen que la venganza es dulce. Y yo lo creo.
Hasta después de que ocurriera este accidente que reveló la causa del "ruido" del violín, nunca supuse que la caja de resonancia del violín pudiera ser tan sensible; sin embargo, tenía a mano suficientes ejemplos similares como para llegar a esa conclusión.
Las cuerdas de piano oxidadas, incluso si solo presentan oxidación localizada, no pueden producir, ni podrán jamás, un sonido puro, libre de ruido. Esto se debe, sin duda, a las irregularidades en el diámetro de las cuerdas causadas por la oxidación irregular. La ley que rige la producción del sonido musical parece explícita y no permite ni siquiera desviaciones mínimas en los agentes que lo producen.
Nuevamente, una lengüeta de órgano ligeramente desafinada se remedia fácilmente. Si su tono es ligeramente demasiado bajo, la eliminación de una cantidad inconmensurable de metal, de
Su extremo libre eleva perceptiblemente su tono; si es demasiado alto, la eliminación de la cantidad justa de metal de su base para aclarar la superficie reduce su tono.
Ahora creo que una caja de resonancia de violín cuidadosamente graduada es igual de sensible que un alambre de piano, una lengüeta de órgano, un tubo de órgano o un diapasón.
"La ciencia nunca ha fabricado un violín."
Señores, por favor, permanezcan en sus asientos.
Como estudiante veterano de violín, la experiencia me obliga a aceptar la veracidad de esta sorprendente afirmación, por muy desagradable que pueda resultar. No soy de los que ignoran la ciencia, pero sí de los que han aprendido a aceptar las afirmaciones científicas con cautela. La ciencia en química, geología, mineralogía, astronomía, botánica, incluso la ciencia en violinología, hasta donde alcanza, me fascina. Hubo un tiempo en que creía en todas las afirmaciones científicas, de cualquier fuente. He vivido hasta la desilusión. Permítanme advertirles del peligro de creer ciegamente en afirmaciones de origen humano. Solo hay una fuente de infalibilidad.
En mi época, la química se enseñaba como una ciencia exacta, tan exacta como las cifras en aritmética. Por lo tanto, cuando leía en los libros de texto de química que el agua es una combinación de dos cuerpos gaseosos elementales, hidrógeno y oxígeno, que el agua equivale a un equivalente de cada uno de estos dos cuerpos,
que por lo tanto el símbolo del agua es H 0, creí que esa afirmación era infaliblemente correcta. Hoy el símbolo del agua ha cambiado; ha cambiado tanto que, sin seguridad, no podría saber 
cómo Ha representa el agua.
En cierta ocasión en la Universidad de Michigan, ocurrió un incidente que me marcó para siempre. Era la hora de la clase de química, en la que tuvimos el honor de contar con la presencia de un anciano que parecía profundamente absorto en el tema que se estaba tratando. Entre otros temas presentados a esa hora estaba el agua destilada, una botella de un galón de la cual reposaba sobre la mesa del profesor. En cuanto terminó la clase, este atento anciano, con una agilidad insospechada, trepó por la barandilla hasta el quirófano, sacó el corcho de la botella de agua destilada, se la olió, negó con la cabeza, agitó la botella, examinó críticamente la "perla" y luego, lentamente...
preguntó: "¿Hay algo en esta agua destilada que sea embriagadora?" Hoy yo soy el anciano. Ustedes son los jóvenes atentos.
De manera pausada pregunto: "¿Es su nuevo H=O más fuerte que el antiguo HO?"
Entonces, el hidrógeno era un elemento, por lo tanto indivisible. Hoy el hidrógeno se separa en argón y boro. No deposites tu fe en la infalibilidad humana; no existe. La ciencia es un Golloso, sin embargo, incluso hoy, la ciencia no puede hacer un violín. Al acercarse al violín en busca de sus secretos, la ciencia, gigante orgulloso, recibe una bofetada; en la cara como recordatorio de que en el mundo hay una cosa que "no le incumbe". En otras direcciones, década tras década, la ciencia avanza a prodigiosos 'límites'; pero en ningún trabajo científico puedo encontrar una solución para
tono de violín.
Tanto los filósofos ingleses como los alemanes descartan el sonido del violín con una evasión latina: "sui generis".
Los filósofos franceses descartan el tono de violín dándole un nombre que pertenece a todos los tonos musicales peculiares,

La ciencia nunca creó un violín. «Doctor, ¿qué creó el violín?» El experimento, por sí solo, creó el violín.
La ciencia llegó más tarde, como suele ocurrir, con una explicación a posteriori de algunos hechos relacionados con el sonido del violín, pero nunca una explicación de todos los hechos relacionados con el sonido del violín.
[A pesar de las declaraciones de autores apologistas, no encuentro ninguna corroboración de que Stradivarius fuera un científico. Que no lo fuera parece estar claramente establecido por su trabajo anterior; que fuera un experimentador incansable, lo demuestra su historia. Que merece reconocimiento lo demuestran algunos de sus violines. Les pido, miembros del Club de Estudiantes de Violín, que indiquen específicamente qué invención o modificación del violín se le debe a Stradivarius.]
¿Silencio? No es de extrañar.
En la lista bibliográfica del violín hay casi 200 libros. Solo de estos libros podemos obtener evidencia. Después de vadear una distancia a través de esta masa de evidencia teñida por escritores que se disculpan por los "trozos" de Strad afirmando que esas anomalías deben haber sido construidas "por contrato", teñida por escritores que poseen una imaginación exuberante, \label{libros} que olvidaron decir que no tienen ningún conocimiento práctico ni de la construcción de violines, ni de la madera para violines, ni del barniz para violines; teñido por escritores que solo conocen la Cremona "amarilla"; de nuevo, por escritores que solo conocen la Cremona "roja"; de nuevo, teñido por promotores comerciales, uno se da cuenta de que ahora sabe menos de los violines Stradivarius que después de haber leído un solo libro.
Ninguno de los dos. Ni Gaurnerius ni Stradivarius parecen haber mejorado el contorno del violín, ni los detalles de construcción, ni el color, ni el barniz. La primera mitad de la obra de Stradivarius se perdió antes de que redujera la curvatura de la tapa armónica a la que Maggini le dio al violín cien años antes. Específicamente, solo puedo determinar que tanto Gaurnerius como Stradivarius redujeron el grosor de la tapa armónica hasta un punto en que las cuerdas más ligeras podían superar la inercia y rigidez de la caja de resonancia; y que, hasta su época, su selección de madera para la caja de resonancia, para violín solista, no tenía parangón. Aparte de estos dos hechos, no encuentro nada más a favor de estos dos renombrados constructores de violines. Si se me pueden mostrar más pruebas a su favor, las reconoceré. Que los constructores de instrumentos musicales italianos, que exigían tres octavas de tono musical de cada cuerda del violín, dieran la pista para la reducción del grosor de la caja de resonancia es sumamente plausible. Que Gaurnerius y Stradivarius fueran los primeros en satisfacer dicha demanda es evidente.
Creo que puedo considerar con seguridad que tanto la reducción del grosor de la mesa como la selección de madera de caja de resonancia de fibra más blanda, eran vistas entonces como
innovaciones; y que, por otros constructores de violines, ambos hombres eran considerados "excéntricos". Pero ambos, como constructores de violines solistas, fueron ganadores, aunque cada uno adoptó un plan diferente para lograr reducciones de grosor en la caja de resonancia. Ole Bull exigió cuatro octavas de tono musical.
de cada cuerda del violín. Encontró satisfacción en un 'José'.
Se dice que Paganni disfrutaba de las cuerdas de violín ligeras. Paganni fue el primer gran solista en explotar el violín Gaurnerius. Honeyman, de forma imprecisa, afirma que la caja de resonancia del Gaurnerius es "más delgada en toda la parte central".
Resulta interesante observar las diversas opiniones de los violinistas sobre las diferentes sonoridades de los violinistas fabricados por los grandes constructores de violines.
Ole Bull se muestra sorprendido al ver a su contemporáneo, DeBeriot, aparecer en público con un Maggini.
Al hablar de su Stradivarius, Ole Bull afirma: «Nunca lo toco en público por su timbre nasal» (véase «Memorias de Ole Bull», de Sarah C., su esposa). Llamo la atención sobre una omisión notable en las declaraciones relativas a los violines Stradivarius: el hablante o escritor suele omitir la mención de a qué «periodo» de Strad pertenece su violín en particular. Solo puedo explicar esta omisión si se parte de la hipótesis de que el hablante o escritor considera que los violines de solo uno de los cuatro «periodos» de Strad poseen una calidad tonal digna de mención.
Caballeros, cometen un error si interpretan que no me uno a la alabanza universalmente aceptada. \label{mejor_violin`}
Admiro a estos dos famosos constructores de violines; pero mis elogios se deben únicamente a una razón concreta: los grandes violines de estos dos grandes fabricantes solo son grandes como violines solistas.
"Pis merece un gran honor, y es más que merecido por su audaz innovación y su correcto juicio sobre la madera de la caja de resonancia,
Mi argumento es el siguiente: los elogios ilimitados que han recibido estos violines solistas han impulsado a un ejército de luthiers a dedicar toda una vida a la reproducción de violines solistas, hasta el punto de que hoy en día apenas hay un buen violín de orquesta en el mercado.
"Doctor, ¿quiere decir que el mejor violín solista no es también el mejor violín de orquesta?" Sí, señor.
"¿Cómo es eso?"
La caja de resonancia del mejor violín solista es demasiado ligera en madera para soportar la tremenda fuerza de
Ondas sonoras de instrumentos de orquesta completa. "Doctor, por favor, explíqueme."
"
Así pues: para que la caja de resonancia del violín produzca tres octavas de tono agradable en cada cuerda, su grosor debe reducirse hasta el punto de debilitarse ante las ondas armónicas de una orquesta completa; y debido a dicha debilidad, sus tonos quedan ahogados por los instrumentos armónicos dominantes.
"Pero, doctor, comprendimos que las ondas sonoras del primer violín se superponen a las ondas de armonía."
Correcto, señor, siempre y cuando las olas de primer violín tengan la fuerza propulsora suficiente para mantenerlas en la cima. El nadador, siempre que su fuerza sea suficiente, puede mantenerse en la cresta de una ola enorme, pero
Cuando la fuerza le falla, su siguiente posición es en el fondo del abismo; luego, en el fondo del mar. Decir que "el violín alcanzó la perfección hace 200 años" significa mucho y poco. Decir que el violín solista alcanzó la perfección hace 200 años es más correcto y no tan engañoso. Deberíamos insistir en especificaciones definidas en este asunto, debido a los diferentes caracteres tonales que se exigen del violín. Solo necesitamos un momento de atención para comprender lo absurdo de tanto esfuerzo mal dirigido para producir violines solistas mientras la demanda de violines de orquesta está en ascenso. Solo unos pocos violines solistas se utilizan realmente en una fecha dada, mientras que el violín de uso general y el violín de orquesta, en grandes cantidades, están en
empleo por horas. "Doctor, por favor, explíqueme qué quiere decir con violín de uso general y violín de orquesta completa." Con mucho gusto, señor.
Así pues: El violín de uso general es aquel cuya afinación en las cuerdas G y D es claramente de carácter grave, mientras que la afinación en las cuerdas A y E es claramente de carácter soprano; el violín de orquesta completa es aquel cuya afinación en las cuerdas G y D es de carácter barítono, mientras que la afinación en las cuerdas A y E es claramente de carácter soprano. "¿Se pueden asignar estas diferentes características tonales a los violines a voluntad? Sí, señor."
"Por favor, díganos cómo."
Al combinar, en el trabajo de construcción de violines, cer-
ciertas reglas que se encuentran en la Lección III, que rigen
tono de los agentes productores de tono.
Gracias a repetidas demostraciones prácticas, puedo asegurar la veracidad de estas reglas. Antes de concluir con la cuestión de la "perfección del violín hace 200 años", permítanme llamar su atención sobre dos de sus nefastas consecuencias.
Primero: Sirve como elemento disuasorio para seguir experimentando.
Segundo: Constituye una atribución ficticia de un valor a unos pocos violines de esa fecha.
Los laboriosos "promotores comerciales" no han perdido oportunidad de confundir al público en lo que respecta a los antiguos valores del violín solista; por lo tanto, el público exige erróneamente que cada solista aparezca solo con
o un 'Joseph*' o un 'Strad' bajo el brazo. Pero,
Hoy hay señales inequívocas de que el público está saliendo de la niebla; algunas personas observadoras, con la valentía de la convicción, pronuncian una "bendición pasajera" sobre el "viejo violín". Mi propia veneración, reverencia, culto, llámenlo como quieran, por el "viejo violín" se vio cruelmente quebrantada hace años al descubrir que el violín puede volverse "demasiado viejo", y también al descubrir que un violín moderno puede lucir "extraordinariamente conservado".
Pero hace unos años el público tuvo la oportunidad de observar la insensatez de seguir exigiendo que el solista utilice únicamente violines antiguos en todas las ocasiones. Fue en el Auditorio de Chicago. Mi amigo, el Sr. Beebe, él mismo un fabricante de violines de renombre, en la primera ocasión, se sentó bastante
cerca del escenario; en una segunda ocasión, se sentó más lejos del escenario; y, en una tercera ocasión, en el asiento más alejado; todo con el propósito de juzgar el tono de un Stradivarius en manos de un solista estrella lo suficientemente sabio como para permitir solo acompañamiento de piano para su viejo violín. El Sr. Beebe se alegra de contar con un oído musicalmente entrenado, y yo, por mi propia experiencia, tengo la certeza de que dicho oído, en presencia del tono del violín, se destaca notablemente. Por lo tanto, cuando el Sr. Beebe afirma que, estando sentado en un asiento alejado del auditorio, ciertos esfuerzos del gran violinista fueron decepcionantes debido a la debilidad en el tono de ese viejo violín, acepto tal afirmación como proveniente de una autoridad confiable. Este incidente muestra el valor ficticio que se le da a los violines antiguos. Ese Stradivarius está valorado en varios miles y Beebe en unos pocos cientos. Yo valoro
La mayoría considera que el violín tiene el tono más satisfactorio.
Señores: Les recomiendo construir violines pensando en su uso específico. Así, al construir un violín para solista, reduzcan el grosor de la caja de resonancia de manera que cada cuerda produzca tres octavas de sonido. Al construir un violín de uso general, reduzcan las dimensiones de la barra y la caja de resonancia debajo de las cuerdas Sol y Re para darles el carácter de bajo, dejando suficiente madera debajo de las cuerdas La y Mi para darles el tono de soprano. Al construir un violín para orquesta, dejen suficiente madera debajo de todas las cuerdas para asegurar el carácter de barítono-soprano.
Después de cincuenta años de trabajo práctico dedicado al violín 

Las peculiaridades del sonido, esas son mis conclusiones en materia de construcción de violines.
Quiero llamar su atención sobre el hecho de que, en violines bien construidos, existe una amplia gama de timbres; y sobre el hecho de que dicha gama se debe a diferencias inherentes en la calidad de la madera de la tapa armónica. Ni usted, ni yo, ni nadie, podemos predeterminar con exactitud la calidad del timbre de una muestra determinada de madera de la tapa armónica. En esto reside la razón por la que la ciencia no puede construir un violín.
Quiero destacar que construir la caja de resonancia "ruidosa" del violín es una tarea sencilla. Para ello, el equipo necesario se detalla en la siguiente lista:
1 gubia de tres cuartos de pulgada. 1 navaja plegable. 1 herramienta de trabajo de veinticinco centavos.
Para esta última "herramienta" debemos recurrir a la importación desde Alemania. Incluso así, resulta un fracaso en este lado del Atlántico. Una vez que esa "herramienta" de veinticinco centavos pisa suelo estadounidense, se convierte en un objeto de dos dólares.
¡Paz para ti, América!
